\documentclass[11pt]{article}
\usepackage[a4paper,margin=26mm]{geometry}
\usepackage{amsmath,amssymb}
\usepackage[hidelinks]{hyperref}
\setlength{\emergencystretch}{3em}
\title{Conjecture 00000001189: its defined support is zero, not $n^2-1$}
\author{lizaixi01, with OpenAI Codex assistance}
\date{23 October 2026}
\begin{document}
\maketitle
\section{The source's own definition and the necessary contradiction}
The source defines minimal support using the minimum of the dimension
of the fixed-point space in the natural $n$-dimensional representation.
It then asserts a generating pair with common fixed-space dimension
zero, each individual fixed space at most one-dimensional, ``giving
support $n^2-1$.'' Both languages use this same preceding definition.

These claims are incompatible already for the legitimate parameters
$n=2$, $q=2$. We prove on the actual group $SL_2(\mathbf F_2)$ that a
generating pair with all those fixed-space properties exists, and that
the actual minimum of fixed dimensions over generating pairs is zero.
It cannot be $2^2-1=3$. More generally, every actual natural fixed
subspace has dimension at most two; the minimum over any nonempty
family of such image fixed spaces therefore also cannot be three.

This refutation uses the source's stated fixed-dimension quantity.
It does not identify it with a different standard support convention,
such as a codimension, an adjoint-representation quantity, or an
algebraic-group dimension. Nor does it assert that Singer-cycle
constructions or the separate generating-pair existence claim fail.
The support-value necessary clause already refutes the conjunction.
The exact original bilingual source is included in the package.

\section{Actual matrices, natural action and fixed spaces}
Lean uses Mathlib's actual special linear group
$G=SL_2(\mathbf F_2)$ and its actual natural action on
$V=\mathbf F_2^2$. For $A\in G$ define the genuine linear submodule
\[
 F_A=\ker(A-I)\le V,\qquad F_{A,B}=F_A\cap F_B.
\]
The membership theorems identify these kernels with the equations
$Av=v$ and $Av=v$, $Bv=v$, respectively. Dimensions are actual
module dimensions over the genuine field \texttt{ZMod 2}.
Since $\dim V=2$, Lean proves both $\dim F_{A,B}\le2$ and
$\dim V^H\le2$ for every actual subgroup $H\le G$.

The subgroup generated by a pair is the actual subgroup closure
$\langle A,B\rangle$. Its fixed space is defined as the intersection
of $F_C$ over all its actual elements $C$. Lean proves
\[
 F_{A,B}=V^{\langle A,B\rangle}.
\]
Indeed, if $v$ is fixed by $A$ and $B$, its actual group-action
stabilizer contains both generators, hence their subgroup closure.
The reverse implication follows because the generators belong to
their closure. This connects common fixed spaces to the source's
fixed space of the generated natural representation image.

\newpage
\section{A genuine generating pair, with all fixed-space conditions}
Take the actual determinant-one transvections
\[
 U=\begin{pmatrix}1&1\\0&1\end{pmatrix},\qquad
 L=\begin{pmatrix}1&0\\1&1\end{pmatrix}
 \quad\text{over }\mathbf F_2.
\]
The formal proof uses Mathlib's special-linear transvection
generation theorem. In dimension two over $\mathbf F_2$, every
transvection is either the identity, $U$, or $L$: the off-diagonal
position has two choices and its coefficient is zero or one.
The subgroup containing $U$ and $L$ therefore contains every
transvection and, by that genuine generation theorem, all of $G$.
Thus their actual subgroup closure is $G$.

For $v=(x,y)$ the actual action gives
\[
 Uv=v\iff y=0,\qquad Lv=v\iff x=0.
\]
Consequently $F_{U,L}=\{0\}$ and its dimension is zero. Each
individual fixed submodule is proper: $(0,1)$ is not fixed by $U$,
and $(1,0)$ is not fixed by $L$. The strict dimension bound for a
proper subspace of $V$ proves that both fixed dimensions are at most
one. The theorem \texttt{actual\_pair\_satisfies\_fixed\_conditions}
collects all these genuine generating-pair and fixed-space properties.

\section{The attained minimum and other image families}
The literal pair-based set of admissible dimensions is
\[
 S=\{\dim F_{A,B}: A,B\in G,\ \langle A,B\rangle=G\}.
\]
The verified pair $U,L$ proves $0\in S$, so this is a nonempty set
of actual dimensions. Nonnegativity gives $\operatorname{IsLeast}(S,0)$.
Lean defines the minimum by $\inf S$ and derives its value zero from
that attained least-element theorem. No default value of an infimum
on an empty set is used.

If the minimum syntax is instead understood as ranging over another
nonempty collection of natural representation images, the conclusion
still cannot be three. Lean separately quantifies over every actual
nonempty set $\mathcal H$ of subgroups of $G$ and proves
\[
 \inf\{\dim V^H:H\in\mathcal H\}\le2.
\]
This uses an actual member of the nonempty collection and the genuine
ambient dimension bound. For the full image $G$ itself, Lean also
proves $V^G=\{0\}$ using the pair-generation bridge above.
Thus neither of these fixed-space interpretations of the source's
minimum can yield $n^2-1$ at $n=2$.

\section{Verification}
Lean 4.33.0 and Mathlib commit
\begin{center}\small\texttt{db584cd6d46c92f209a44c0f1c829460d327499d}\end{center}
are pinned. The package includes a new root build, transitive axiom
audits for every theorem and full official
\texttt{leanchecker --fresh Main} replay. No \texttt{sorry},
\texttt{native\_decide} or extra axiom is used. OpenAI Codex assisted;
a separate supervising Codex agent reviewed source, report and PDF.
This is AI review, not human certification. Organizer acceptance is pending.
\end{document}
